\section{Trace the first byte before processing the whole message}
The ASCII byte for the character 1 is hexadecimal \texttt{31}. Starting from
\texttt{ffffffff}, XOR produces \texttt{ffffffce}. Its low bit is zero, so the
first shift produces \texttt{7fffffe7} without feedback. The new low bit is one:
the next shift gives \texttt{3ffffff3}, and XOR with the feedback mask gives
\texttt{d2477cd3}. Continue inspecting the \emph{current} low bit at each step.

\begin{table}[htbp]
\centering
\caption{Eight reflected steps after XORing the first byte, \texttt{31}.}
\tagpdfsetup{table/header-rows={1}}
\begin{tabular}{rl}
\toprule
Completed bit steps & Raw state\\
\midrule
0 & \texttt{ffffffce}\\
1 & \texttt{7fffffe7}\\
2 & \texttt{d2477cd3}\\
3 & \texttt{849b3d49}\\
4 & \texttt{aff51d84}\\
5 & \texttt{57fa8ec2}\\
6 & \texttt{2bfd4761}\\
7 & \texttt{f8462090}\\
8 & \texttt{7c231048}\\
\bottomrule
\end{tabular}
\end{table}

For the complete nine-byte ASCII sequence \texttt{123456789}, the raw result is
\texttt{340bc6d9}; complementing it produces \texttt{cbf43926}. These values are
separate observations at the update and message boundaries. If a program reports the
raw result as its final checksum, it can have a correct inner loop and an incorrect
file result.

The input contains nine bytes, without a trailing newline or terminating zero. Adding
either byte changes the message being checked. When comparing a result with another
program, first make sure the byte sequences agree, then check the CRC parameters.
A matching filename or visible text does not establish equal file contents.
